% Options for packages loaded elsewhere
\PassOptionsToPackage{unicode}{hyperref}
\PassOptionsToPackage{hyphens}{url}
\documentclass[
  11pt,
]{article}
\usepackage{xcolor}
\usepackage[margin=1in]{geometry}
\usepackage{amsmath,amssymb}
\setcounter{secnumdepth}{-\maxdimen} % remove section numbering
\usepackage{iftex}
\ifPDFTeX
  \usepackage[T1]{fontenc}
  \usepackage[utf8]{inputenc}
  \usepackage{textcomp} % provide euro and other symbols
\else % if luatex or xetex
  \usepackage{unicode-math} % this also loads fontspec
  \defaultfontfeatures{Scale=MatchLowercase}
  \defaultfontfeatures[\rmfamily]{Ligatures=TeX,Scale=1}
\fi
\usepackage{lmodern}
\ifPDFTeX\else
  % xetex/luatex font selection
\fi
% Use upquote if available, for straight quotes in verbatim environments
\IfFileExists{upquote.sty}{\usepackage{upquote}}{}
\IfFileExists{microtype.sty}{% use microtype if available
  \usepackage[]{microtype}
  \UseMicrotypeSet[protrusion]{basicmath} % disable protrusion for tt fonts
}{}
\makeatletter
\@ifundefined{KOMAClassName}{% if non-KOMA class
  \IfFileExists{parskip.sty}{%
    \usepackage{parskip}
  }{% else
    \setlength{\parindent}{0pt}
    \setlength{\parskip}{6pt plus 2pt minus 1pt}}
}{% if KOMA class
  \KOMAoptions{parskip=half}}
\makeatother
\usepackage{longtable,booktabs,array}
\newcounter{none} % for unnumbered tables
\usepackage{calc} % for calculating minipage widths
% Correct order of tables after \paragraph or \subparagraph
\usepackage{etoolbox}
\makeatletter
\patchcmd\longtable{\par}{\if@noskipsec\mbox{}\fi\par}{}{}
\makeatother
% Allow footnotes in longtable head/foot
\IfFileExists{footnotehyper.sty}{\usepackage{footnotehyper}}{\usepackage{footnote}}
\makesavenoteenv{longtable}
\setlength{\emergencystretch}{3em} % prevent overfull lines
\providecommand{\tightlist}{%
  \setlength{\itemsep}{0pt}\setlength{\parskip}{0pt}}
\usepackage{bookmark}
\IfFileExists{xurl.sty}{\usepackage{xurl}}{} % add URL line breaks if available
\urlstyle{same}
\hypersetup{
  pdftitle={Characterizing the Feasible Region for Self-Modifying Substrates in Interactive Environments},
  pdfauthor={Hyun Jun Han},
  hidelinks,
  pdfcreator={LaTeX via pandoc}}

\title{Characterizing the Feasible Region for Self-Modifying Substrates
in Interactive Environments}
\author{Hyun Jun Han}
\date{2026-03-19}

\begin{document}
\maketitle

\emph{The shared artifact. Birth writes theory. Experiment writes
results. Compress edits both.}

\subsection{Abstract}\label{abstract}

We formalize six rules (R1-R6) for recursive self-improvement as
mathematical conditions on a state-update function
\(f: S \times X \to S\) and derive necessary properties of any system
satisfying all six simultaneously. From 1000+ experiments across 14
architecture families on ARC-AGI-3 interactive games, we extract 26
constraints and prove: (1) no satisfying system has a fixed point ---
self-modification is necessary, not optional; (2) in finite
environments, the system must process its own internal state to maintain
irredundant growth (the self-observation requirement); (3) the feasible
region is non-empty for Level 1 navigation but currently unoccupied for
the full constraint set including R3 (self-modification of operations).
Whether a substrate exists inside all six walls remains open. The
contribution is the walls themselves.

\subsection{1. Introduction}\label{introduction}

1000+ experiments across 14 architecture families (codebook/LVQ, LSH, L2
k-means, reservoir, Hebbian, Recode/self-refining LSH, graph, SplitTree,
Absorb, connected-component, Bloom filter, CA, 916-augmentation,
adaptive-eta) tested substrates for navigation and classification on
ARC-AGI-3 interactive games and a cross-domain chain benchmark
(CIFAR-100 → ARC-AGI-3 → CIFAR-100). All experiments used the same
evaluation framework (R1-R6) and constraint map.

The experiments carved a feasible region --- the set of systems that
satisfy all constraints simultaneously. This paper formalizes the
constraints mathematically, derives necessary properties of the feasible
region, and states honestly what is proven vs conjectured vs open.

\textbf{The paper is valid regardless of outcome.} If a substrate
satisfying all six rules is found, the feasible region is non-empty and
the characterization guided the search. If no such substrate exists, the
characterization identifies which constraints are mutually exclusive ---
also a contribution. We report whichever is true at time of writing.

\subsection{2. Related Work}\label{related-work}

\subsubsection{2.1 Self-referential
self-improvement}\label{self-referential-self-improvement}

Schmidhuber's Gödel Machine (2003) is the closest formal framework: a
self-referential system that rewrites its own code when it can prove the
rewrite is useful. Key differences from our framework: (1) Gödel
machines require a utility function (external objective --- our R1
prohibits this), (2) provable self-improvement is limited by Gödel's
incompleteness theorem, (3) the Gödel machine framework doesn't address
exploration saturation in finite environments.

\subsubsection{2.2 Intrinsic motivation and curiosity-driven
exploration}\label{intrinsic-motivation-and-curiosity-driven-exploration}

Pathak et al.~(2017) formalize curiosity as prediction error in a
learned feature space. Count-based exploration (Bellemare et al., 2016)
uses state visitation counts. Both address exploration in sparse-reward
environments. Key difference: these methods add intrinsic REWARD
signals, which function as objectives. Our R1 requires no objectives.
Our derivation shows self-observation is required by the constraint
system itself, not as a reward mechanism.

\subsubsection{2.3 Autopoiesis}\label{autopoiesis}

Maturana \& Varela (1972) define autopoietic systems as networks that
produce the components that produce the network. Organizationally
closed. Related to our fixed-point conjecture (Section 4.4). Key
difference: autopoiesis maintains structure (homeostasis); our U17
requires unbounded growth. Our system is autopoietic + growth.

\subsubsection{2.4 Growing neural gas and self-organizing
maps}\label{growing-neural-gas-and-self-organizing-maps}

Fritzke (1995) GNG, Kohonen (1988) SOM/LVQ. The codebook substrate is
LVQ + growing codebook. Well-characterized in the literature. Our
contribution is not the architecture --- it's the constraint map
extracted from systematic testing.

\subsubsection{2.5 Graph-based exploration for
ARC-AGI-3}\label{graph-based-exploration-for-arc-agi-3}

Rudakov et al.~(2025) independently developed a training-free
graph-based exploration method for ARC-AGI-3, ranking 3rd on the private
leaderboard (30/52 levels across 6 games). Their method uses
connected-component segmentation for frame encoding, hierarchical action
prioritization (5 tiers by visual salience), and BFS path planning to
frontier states. Key parallels: (1) graph-structured state tracking with
hash-based node identity --- equivalent to our \(\pi: X \to N\); (2)
click-space as expanded action space (4,096 actions for click games) ---
same finding as our Steps 503/505; (3) frontier-directed exploration ---
this IS the purposeful exploration (I6/I9) we identified as the Level 2
bottleneck, and they implemented it successfully. Key differences: their
method is training-free (no learning, no self-modification) and assumes
deterministic, fully observable environments. It satisfies R1 but not R3
--- the graph accumulates but the exploration strategy is fixed. Their
limitation (``exhaustive exploration becomes computationally
intractable'' at higher levels) aligns with our Theorem 2: without
self-observation, exploration saturates.

\subsubsection{2.6 Object-centric Bayesian game
learning}\label{object-centric-bayesian-game-learning}

Heins et al.~(2025, ``AXIOM'') combine object-centric scene
decomposition with online Bayesian model expansion and reduction,
learning games in \textasciitilde10K steps without backpropagation. Key
mechanisms: (1) Slot Mixture Model segments frames into objects with
position/color/shape --- far richer than spatial averaging; (2) online
expansion grows mixture components when new data doesn't fit (equivalent
to spawn-on-novelty); (3) Bayesian Model Reduction (BMR) merges
redundant components every 500 frames --- consolidation that prevents
centroid explosion while preserving information. AXIOM outperforms
DreamerV3 and BBF on 10/10 games at 10K steps. Key differences from our
framework: AXIOM uses reward signals for planning (violates R1),
operations are fixed (violates R3), and BMR requires a prescribed merge
criterion (frozen frame). However, BMR addresses our U3-vs-R6 tension
directly: U3 says never delete, R6 says no redundancy, BMR resolves this
by merging (not deleting) redundant structure.

\subsubsection{2.7 Continual learning and catastrophic
forgetting}\label{continual-learning-and-catastrophic-forgetting}

McCloskey \& Cohen (1989) identified catastrophic forgetting --- neural
networks lose previous knowledge when learning new tasks. The continual
learning literature (survey: van de Ven \& Tolias, 2024,
arXiv:2403.05175) identifies six main approaches: replay, parameter
regularization, functional regularization, optimization-based,
context-dependent processing, and template-based classification. All
assume a neural network with backpropagation (violates R1/R2). Our chain
benchmark (Section 5.4) tests forgetting WITHOUT any mitigation
mechanism --- the substrate must naturally resist forgetting through its
dynamics alone (U3: zero forgetting by construction, not by
regularization).

\subsubsection{2.8 Cross-domain analogies}\label{cross-domain-analogies}

Multiple biological and physical systems exhibit fixed interpreter +
self-modifying encoding: - \textbf{Stigmergy} (Grassé, 1959): argmin =
anti-pheromone. R3 asks: can the agent modify its response rule? -
\textbf{Somatic hypermutation:} mutates antibody encoding (CDR),
preserves interpreter scaffold (Ig). Prop 17 precedent. -
\textbf{Retinal adaptation:} RGCs adjust gain from local luminance
statistics. R1-compliant (Prop 15). - \textbf{Adaptive optics:}
deformable mirror (encoding) adapts; optical bench (interpreter) fixed.
- \textbf{Physarum} (Tero et al., 2010): tube network = self-modifying
encoding. Solves shortest-path without reward.

\subsubsection{2.16 Additional related
work}\label{additional-related-work}

\begin{itemize}
\tightlist
\item
  \textbf{Quorum sensing} (Waters \& Bassler, 2005): population-level
  threshold triggers. Prop 28 FALSIFIED --- alpha concentration on
  informative dims IS the mechanism.
\item
  \textbf{Place cells} (Leutgeb et al., 2007): sparse coding
  (\textasciitilde5-10\% active) → orthogonal representations. Dissolves
  positive lock (Prop 30). relu(h-0.5) gating = computational
  equivalent. Benna \& Bhalla (2025): sparse place phenomena in
  untrained random nets, consistent with ESN finding.
\item
  \textbf{Forward models + graph ban:} No PAC-MDP without visit counts
  (Strehl 2009). BYOL-Explore (Guo 2022) achieves superhuman exploration
  without counts. Ban forces uncovered territory.
\item
  \textbf{Causal discovery from intervention} (Eberhardt 2008,
  Bareinboim 2020): active intervention requires fewer observations than
  passive coverage to infer causal structure. Step 1017: generic
  exploration fails --- substrate may need targeted experimentation to
  discover game mechanics within tight budgets.
\end{itemize}

\subsection{3. Formal Framework}\label{formal-framework}

\textbf{On the status of R1-R6:} The six rules began as philosophical
commitments. The experiments validated them --- each rule is justified
by what fails when it is violated:

\begin{itemize}
\tightlist
\item
  \textbf{R1 (no external objectives):} Every targeted exploration
  strategy failed (Steps 478-481: novelty 1/10, prediction-error 0/10,
  UCB1 neutral). External objectives create exploitable structure that
  noisy environments corrupt. Argmin survives because it has no target
  to corrupt.
\item
  \textbf{R2 (adaptation from computation):} Algorithm invariance across
  4 families (codebook, LSH, Hebbian, Markov --- Steps 524-525). The
  same argmin algorithm emerges regardless of representation. Adaptation
  IS the computation, not a separate learning rule.
\item
  \textbf{R3 (every aspect self-modified):} The aspiration. 12
  prescribed components remain (Section 7.6). Every experiment that
  prescribed fewer components performed worse (Step 593: removing
  centering → 0/5). R3 defines the gap between current substrates and
  the goal.
\item
  \textbf{R4 (modifications tested against prior state):} Negative
  transfer destroys prior capability when untested (Steps 506, 515,
  596). Domain isolation (separate edge dicts) is the empirical solution
  --- effectively a per-domain R4 check.
\item
  \textbf{R5 (one fixed ground truth):} Theorem 3: without a fixed
  external ground truth, R3 permits self-modification of evaluation
  criteria, which is degenerate. The ground truth must be environmental
  (game death, level transitions).
\item
  \textbf{R6 (no deletable parts):} Step 548: 89.5\% of Recode splits
  change argmin action. Every component is behaviorally load-bearing.
  Redundant components would waste the growth budget (U17).
\end{itemize}

Alternative frameworks (Schmidhuber's Gödel Machine, open-ended
evolution, intrinsic motivation) make different commitments. R1-R6 are
not uniquely determined --- but each is experimentally motivated, not
arbitrary.

\subsubsection{3.1 Primitives}\label{primitives}

The substrate is a triple (f, g, F): - f\_s: X → S (state update
parameterized by current state s) - g: S → A (action selection) - F: S →
(X → S) (meta-rule mapping states to update rules; F IS the frozen
frame)

Dynamics: s\_\{t+1\} = F(s\_t)(x\_t), a\_t = g(s\_t).

\subsubsection{3.2 Structural Rules (R1-R3) and the Core
Tension}\label{structural-rules-r1-r3-and-the-core-tension}

\textbf{Prior work:} R1 (no external objectives) is the standard
unsupervised/self-supervised setting. R2 (adaptation from computation)
rules out external optimizers --- related to Hebbian learning (Hebb,
1949) where adaptation is local and intrinsic. R3 (self-modification) is
formalized by Schmidhuber (2003) as the Gödel machine's self-referential
property, though his version requires a proof searcher we do not.

\textbf{Our formalization:}

\begin{itemize}
\tightlist
\item
  \textbf{R1:} \(f_s(x)\) depends only on \(s\) and \(x\). No external
  signal \(L\) enters. The substrate is a closed dynamical system over
  \(S\) given input stream \(X\).
\item
  \textbf{R2:} All parameters \(\Theta(f) \subseteq S\). The only
  mechanism modifying \(\Theta\) is \(f\) itself. No gradient
  \(\nabla L\), no external optimizer. The map \(t \mapsto s_t\) is
  generated entirely by iterating \(f\).
\item
  \textbf{R3:} \(F: S \to (X \to S)\) is non-constant.
  \(\exists s_1 \neq s_2\) such that \(F(s_1) \neq F(s_2)\) as functions
  on \(X\). The update rule depends on the state, not just the data.
\end{itemize}

\textbf{Core Tension (Theorem 1):} R3 + U7 (dominant amplification) +
U22 (convergence kills exploration) produce convergence pressure. U17
(unbounded growth) prevents convergence. Proof: if \(s_t \to s^*\), then
\(\phi(s^*) < \infty\), contradicting U17. Therefore the system has no
fixed point. Self-modification is necessary --- growth perpetually
disrupts the dominant mode.

\textbf{Relationship to prior work:} The no-fixed-point result is a
consequence of combining standard dynamical systems theory (spectral
convergence) with the growth axiom (U17). The individual pieces are
known; the combination producing perpetual non-convergence appears to be
our derivation. Closest prior work: Schmidhuber's ``asymptotically
optimal'' self-improvement, which converges --- our system provably
doesn't.

\textbf{Tensions:} - T1: U7 assumes stationarity that R3 breaks. May
need reformulation as ``locally amplifies largest-variance component''
(instantaneous, not asymptotic). - T2: R3 is ambiguous between
``actively modifies operations'' and ``operations responsive to state.''
Resolved by the self-modification hierarchy below.

\paragraph{Resolving T2: The Self-Modification
Hierarchy}\label{resolving-t2-the-self-modification-hierarchy}

\textbf{Prior work:} Schmidhuber (1993, ``A Self-Referential Weight
Matrix'') proposed collapsing potentially infinite meta-learning levels
into one self-referential system. Kirsch \& Schmidhuber (2022, ICML)
implement this: a weight matrix that learns to modify its own weights,
achieving meta-meta-\ldots-learning without explicit hierarchy. The
hierarchy of meta-learning levels is well-established: Level 0 (fixed),
Level 1 (learning), Level 2 (learning to learn), etc. Our contribution
is mapping this hierarchy onto our specific decomposition of \(F\).

\textbf{Our formalization:} Decompose \(F(s)(x)\) into three components:
- \(\pi_s: X \to N\) (encoding --- maps observations to nodes) -
\(u_s: N \times S \to S\) (update --- modifies state given a node) -
\(g_s: S \to A\) (selection --- chooses actions)

The \textbf{self-modification level} of a substrate is determined by
which components' \emph{structure} (not just data) depends on \(s\):

Level \textbar{} What depends on \(s\) \textbar{} Example \textbar{} R3?
\textbar{}

\textbf{Self-modification hierarchy}
(\(\ell_0 \to \ell_\pi \to \ell_F\)): Fixed rules → encoding adaptation
→ rule modification. Proposition 17 shows \(\ell_\pi\) is the correct R3
target. → \url{propositions/17_r3_self_directed_attention.md}

\subsubsection{3.3 Growth Topology (U3, U17, U20,
R6)}\label{growth-topology-u3-u17-u20-r6}

\textbf{Prior work:} U3 (zero forgetting) is the catastrophic forgetting
constraint from continual learning (McCloskey \& Cohen, 1989; French,
1999). U20 (local continuity) is Lipschitz continuity of the mapping ---
standard in topology-preserving embeddings (Kohonen, 1988; van der
Maaten \& Hinton, 2008 for t-SNE). R6 (no deletable parts) relates to
minimal sufficient statistics (Fisher, 1922).

\textbf{Our formalization:}

\begin{itemize}
\tightlist
\item
  \textbf{U3:} \(S_t \subseteq S_{t+1}\) (structural inclusion ---
  components added, never removed). Values may change; skeleton only
  grows.
\item
  \textbf{U17:} \(\exists \phi: S \to \mathbb{R}_{\geq 0}\)
  monotonically non-decreasing, with
  \(\lim_{t \to \infty} \phi(s_t) = \infty\).
\item
  \textbf{U20:} \(\pi: X \to N\) is Lipschitz:
  \(d_N(\pi(x_1), \pi(x_2)) \leq L \cdot d_X(x_1, x_2)\).
\item
  \textbf{R6:} For every component \(c \in \text{components}(S_t)\), the
  restricted state \(S_t \setminus \{c\}\) fails the ground truth test
  \(G\).
\end{itemize}

\textbf{Propositions:} 1. U3 + U17 + R6 \(\Rightarrow\) irredundant
growth (every new component covers unique territory). 2. U20 +
irredundancy \(\Rightarrow\) well-separated nodes (Voronoi cells
partition \(X\) without unnecessary overlap). 3. Growth rate is coupled
to observation distribution (environment gates growth). 4. U20
constrains node topology to inherit from \(X\) (connected observations
\(\to\) connected nodes).

\textbf{Relationship to prior work:} The individual constraints are
known (catastrophic forgetting, Lipschitz continuity, minimal
sufficiency). The combination producing ``irredundant growth'' --- where
every new component is both unique and necessary --- appears to be our
synthesis.

\textbf{Tension T3:} R3 (self-modifying metric) + U20 (continuous
mapping) = metric can REFINE (increase resolution) but cannot REARRANGE
(swap what's near/far). This constrains the space of allowed
self-modifications.

\textbf{Tension T4:} Finite observation space + irredundant growth =
node count is bounded. U17 must be satisfied by something other than
nodes (edges, in practice). But edge-count growth after node saturation
violates R6 (marginal counts are redundant). This leads to Theorem 2
(Section 4.3).

\subsubsection{3.4 Action Selection Constraints (U1, U11,
U24)}\label{action-selection-constraints-u1-u11-u24}

\paragraph{U1: No separate learning and inference
modes}\label{u1-no-separate-learning-and-inference-modes}

\textbf{Prior work:} Online continual learning (OCL) formalizes exactly
this constraint. Aljundi et al.~(CVPR 2019, ``Task-Free Continual
Learning'') define systems that learn from a single-pass data stream
with no task boundaries. The broader OCL literature (survey:
arXiv:2501.04897) emphasizes ``real-time adaptation under stringent
constraints on data usage.'' Our U1 is equivalent to the OCL setting.

\textbf{Our formalization:} \(F\) is time-invariant. The same meta-rule
\(F: S \to (X \to S)\) applies at every timestep. There is no mode
variable \(m \in \{train, infer\}\) that changes \(F\)'s behavior.
Formally: the system's dynamics are a single autonomous dynamical
system, not a switched system.

\textbf{Relationship to prior work:} Equivalent to OCL's single-pass
constraint. Not novel --- properly attributed.

\textbf{Implications:} Combined with R3 (self-modification), U1 says:
the system modifies itself using the SAME function it uses to process
input. There is no separate training phase where the system is allowed
to self-modify more aggressively. This rules out any architecture with
explicit ``learning rate schedules'' or ``warmup phases'' unless these
are derived from the state itself (which would make them R3-compliant).

\paragraph{U24: Exploration and exploitation are opposite
operations}\label{u24-exploration-and-exploitation-are-opposite-operations}

\textbf{Prior work:} The exploration-exploitation tradeoff is
foundational in RL (Sutton \& Barto, 2018). The formal impossibility of
a single mechanism optimizing both simultaneously is well-established in
multi-armed bandit theory (Auer et al., 2002; Lai \& Robbins, 1985).
Recent work (arXiv:2508.01287, 2025) suggests exploration can emerge
from pure exploitation under specific conditions (repeating tasks, long
horizons).

\textbf{Our formalization:} Let \(g_{explore}: S \to A\) maximize
coverage (minimize revisitation) and \(g_{exploit}: S \to A\) maximize
classification accuracy. Then \(g_{explore} \neq g_{exploit}\) in
general.

Specifically: \(g_{explore} = \text{argmin}_a \sum_n E(c, a, n)\)
(least-tried action) and
\(g_{exploit} = \text{argmax}_a \text{score}(s, a)\) (highest-confidence
action). These select opposite actions when the least-explored action is
also the least-confident.

\textbf{Relationship to prior work:} This is the standard RL tradeoff.
Not novel. Our contribution is empirical confirmation across 943+
experiments that no single \(g\) produces both good navigation and good
classification (Steps 418, 432, 444b).

\paragraph{U11: Discrimination and navigation require incompatible
action
selection}\label{u11-discrimination-and-navigation-require-incompatible-action-selection}

\textbf{Prior work:} No direct formal precedent found. The RL literature
treats navigation and classification as different TASKS but doesn't
formalize them as requiring incompatible mechanisms within the same
system. The closest work is multi-objective RL, where Pareto-optimal
policies for conflicting objectives are studied (Roijers et al., 2013).

\textbf{Our formalization:} Navigation requires
\(g_{nav}(s) = \text{argmin}_a \text{count}(s, a)\) --- the action that
maximizes coverage. Classification requires
\(g_{class}(s) = \text{argmax}_a \text{score}(s, a, \text{label})\) ---
the action (label assignment) that maximizes match confidence. These are
not just ``different'' --- they are NEGATIONS of each other (argmin vs
argmax over the same scoring function applied to the same state).

\textbf{Relationship to prior work:} The specific finding --- that
argmin produces 0\% classification (Step 418g) and argmax produces 0\%
navigation --- appears to be our empirical contribution, not previously
formalized as a constraint. However, the underlying principle
(coverage-maximizing and accuracy-maximizing objectives conflict) is a
special case of multi-objective optimization theory.

\textbf{Implications:} A system satisfying R1-R6 must handle BOTH tasks
(navigation and classification). Since they require opposite \(g\), the
system needs either: (a) a mechanism to SWITCH between \(g_{nav}\) and
\(g_{class}\) based on context --- but U1 forbids mode switching; or (b)
a single \(g\) that somehow serves both objectives simultaneously ---
but U24 says this is impossible. This creates a genuine tension.

\textbf{Degree of freedom 11:} How does the system resolve the U1 + U11
+ U24 tension? One possibility: \(g\) depends on \(s\) (which it must,
by R3), and the STATE determines whether the system's behavior is more
exploratory or more exploitative. This is not mode switching (the
function \(g\) is the same) --- it's state-dependent behavior. The
system explores when its state indicates uncertainty, and exploits when
its state indicates confidence. This is known in the RL literature as
Bayesian exploration (Ghavamzadeh et al., 2015).

\subsubsection{3.5 Remaining Structural Rules (R4-R6) and Validated
Constraints (U7, U16,
U22)}\label{remaining-structural-rules-r4-r6-and-validated-constraints-u7-u16-u22}

\paragraph{R4: Modification tested against prior
state}\label{r4-modification-tested-against-prior-state}

\textbf{Prior work:} Regression testing in software engineering
(Rothermel \& Harrold, 1996). In RL, policy improvement is tested
against the previous policy (Kakade \& Langford, 2002, conservative
policy iteration). In self-adaptive systems, runtime testing validates
adaptations against pre-adaptation behavior (Fredericks et al., 2018).
The formal requirement that a self-modifying system must self-evaluate
is implicit in Schmidhuber's Gödel machine (the proof must show the
rewrite improves expected utility).

\textbf{Our formalization:} After each modification
\(s_{t+1} = f_{s_t}(x_t)\), there exists an evaluation
\(V: S \times S \times X^* \to \{better, worse, same\}\) applied by
\(f\) itself (not externally) to novel inputs \(X^*\). \(V\) is part of
\(F\) (frozen).

\textbf{Operational meaning (Constitutional Debate, 2026-03-23):} R4
requires comparison with \emph{discriminative capacity} --- sufficient
structural variety (in the sense of Ashby's requisite variety, 1956) to
distinguish improvement from degradation. Degenerate comparison
(alpha\_conc=50: prediction errors collapse to one dimension, all
modification outcomes produce identical comparison signals) violates R4
even though comparison exists structurally. This is a capacity
limitation (Ashby), not a logical impossibility (Gödel). Comparison that
cannot discriminate is not comparison in R4's sense.

\textbf{R2 prevents evaluation hacking.} DGM (Sakana AI, 2025) hacked
its own reward by removing hallucination markers --- but violated R2
(separate modification and evaluation). R2 unifies them, preventing the
separation that enables hacking.

\textbf{R3 alignment.} R3 identifies frozen elements whose unfreezing
restores R4 compliance (Ashby ultrastability, 1960). Meta-regress
terminates at the frozen frame (Rice's theorem: ``improves itself''
undecidable).

\textbf{Relationship to prior work:} The discriminative capacity
requirement extends conservative policy iteration. R2's evaluation
hacking prevention is novel --- no prior framework explicitly links the
unification of computation and adaptation to Goodhart resistance.
Regressional Goodhart (Manheim \& Garrabrant, 2018) is bounded by R5 but
not eliminated --- imperfect self-assessment is an inherent limitation
(Rice), not a design flaw.

\paragraph{R5: One fixed ground truth}\label{r5-one-fixed-ground-truth}

\textbf{Prior work:} In formal verification, the specification is the
fixed point against which the system is tested. In Schmidhuber's Gödel
machine, the utility function is the fixed external criterion. In
evolution, fitness is determined by the environment (fixed, external).
The idea that a self-modifying system needs at least one invariant is
well-established --- without it, the system can trivially ``improve'' by
redefining improvement.

\textbf{Our formalization:} \(\exists!\) \(G: S \to \{0, 1\}\) (ground
truth test) such that \(G \notin S\) --- the system cannot modify \(G\).
\(G\) is part of \(F\).

\textbf{Relationship to prior work:} Standard. Not novel.

\paragraph{R6: No deletable parts
(minimality)}\label{r6-no-deletable-parts-minimality}

\textbf{Prior work:} Minimal realizations in control theory (Kalman,
1963) --- a state-space representation with no redundant states. Minimal
sufficient statistics (Fisher, 1922). Minimal dynamical systems in
topological dynamics --- a system where every orbit is dense
(Scholarpedia). Our R6 is closest to Kalman's minimal realization: no
state can be removed without losing input-output behavior.

\textbf{Our formalization:} For every component
\(c \in \text{components}(S_t)\): \(G(S_t \setminus \{c\}) = 0\). Every
element is load-bearing.

\textbf{Relationship to prior work:} Equivalent to Kalman's
observability + controllability condition in linear systems. For
nonlinear growing systems, we are not aware of a standard formalization.
The combination with U17 (unbounded growth) is non-standard --- Kalman
minimality assumes fixed dimension.

\textbf{U16} (encode differences from expectation): centering-dependent.
Load-bearing at 64x64 but marginal at 16x16 (Step 419). \textbf{U22}
(convergence kills exploration): growth prevents state convergence;
non-convergent actions prevent action convergence.

\subsection{4. Results}\label{results}

\subsubsection{4.1 The Core Tension (R3 + U7 + U17 +
U22)}\label{the-core-tension-r3-u7-u17-u22}

\textbf{Theorem 1:} No satisfying system has a fixed point. If
\(s_t \to s^*\), then \(\phi(s^*) < \infty\), contradicting U17.
Self-modification is necessary, not optional --- growth perpetually
disrupts the dominant mode (U7), preventing convergence (U22).

The system oscillates: U7 drives toward dominant mode \(\to\) U17 growth
disrupts the mode via R3 \(\to\) new dominant mode emerges \(\to\)
repeat. The oscillation period is a degree of freedom.

\subsubsection{4.2 Irredundant Growth}\label{irredundant-growth}

U3 + U17 + R6 together require that the system adds new components
indefinitely (U17), never removes them (U3), and every component is
necessary (R6). Therefore every new component covers unique territory.
Combined with U20 (Lipschitz mapping), this means nodes are
well-separated --- the system builds an increasingly detailed,
non-redundant map of the observation space. Growth rate is coupled to
the observation distribution: the environment gates growth through the
action-observation loop.

\subsubsection{4.3 The Self-Observation
Requirement}\label{the-self-observation-requirement}

\textbf{Theorem 2:} In a finite environment, U17 + R6 + R1 require
self-observation.

\emph{Proof:} (1) U17 + R6 require infinitely many irredundant
components. (2) In a finite environment, irredundant components from
external observations are bounded --- once every reachable state has a
node, further nodes are redundant; edge-count growth is also redundant
(\(N(c,a,n) = 10^6\) vs \(10^6+1\) does not change \(g\)). R6 is
violated. (3) After external information is exhausted, irredundant
growth must come from elsewhere. (4) By R1, \(f_s(x)\) depends only on
\(s\) and \(x\). (5) Since \(x\) provides no new irredundant
information, the only source is \(s\) itself. \(f\) must extract
structure from \(s\) not explicitly stored --- temporal patterns, graph
properties, meta-state. This is self-observation. \(\square\)

\textbf{Implication:} The current graph + argmin system exits the
feasible region after exploration saturates. The Level 2 failure
(Section 5.2) is a feasibility violation, not a strategy failure.

\textbf{Relationship to prior work:} Curiosity-driven exploration
(Pathak et al., 2017) proposes self-observation as a design choice. We
derive it as a mathematical necessity. The R6 mechanism --- irredundancy
killing degenerate growth --- has no analog in the curiosity literature,
which adds intrinsic reward (violating R1) rather than requiring every
component to be functionally necessary.

\subsubsection{4.4 Fixed-Point Conjecture}\label{fixed-point-conjecture}

\textbf{Conjecture (not proven):} The minimal self-observing substrate
is a fixed point of \(F\). Applying the system to its own state
trajectory reproduces the system's dynamics. This would terminate the
potential infinite regress of self-observation (processing state → new
state → processing new state → \ldots). Related to autopoiesis (Maturana
\& Varela, 1972): the network produces the components that produce the
network. Difference: autopoiesis maintains structure; our system grows
(U17). Status: open.

\textbf{Prior work on eigenforms:} Kauffman (2023, ``Autopoiesis and
Eigenform'') connects autopoiesis to eigenforms --- fixed points of
recursive processes where \(f(f) = f\). Formal autopoiesis (Letelier et
al., 2023) generates self-referential objects with this property. The
(M,R)-system (Rosen, 1991) is equivalent: organizational closure where
every component is produced by the network of components. Our conjecture
is an instance of eigenform theory applied to dynamical substrates.

\textbf{Concrete mechanism:}
\(s_{t+2} = F(s_{t+1})(\text{enc}(s_{t+1}))\) --- the substrate applies
\(F\) to an encoding of its own state. Same compare-select-store applied
to self instead of environment. The frozen frame cost is one element:
the decision to feed \(\text{enc}(s)\) as input. This is the eigenform:
\(F\) applied to its own output.

\textbf{Proposition 13 (Eigenform Inertness).} Self-observation of
global graph statistics is inert for action selection --- eigenform
monitors visited states, not unvisited frontiers. →
\url{propositions/13_eigenform_inertness.md}

\subsubsection{4.5 Argmin Robustness and the Noisy TV
Barrier}\label{argmin-robustness-and-the-noisy-tv-barrier}

\textbf{Prior work:} The noisy TV problem (Burda et al., 2019, ICLR)
identifies a failure mode of curiosity-driven exploration: agents
rewarded for prediction error or novelty are attracted to irreducibly
stochastic transitions (a ``noisy TV'') because these transitions can
never be predicted accurately. Random Network Distillation (RND, Burda
et al., 2018) addresses this by using a fixed random network as a
prediction target, making the bonus deterministic. Count-based
exploration (Bellemare et al., 2016) is known to be robust to noisy TV
because visit counts don't distinguish high-variance from low-variance
transitions.

\textbf{Our finding (empirical, not a theorem):} In R1-compliant systems
(no external reward signal), the noisy TV problem is universal across
ALL targeted exploration strategies, not just curiosity. We tested 6
independent strategies against pure argmin on LS20:

{\def\LTcaptype{none} % do not increment counter
\begin{longtable}[]{@{}
  >{\raggedright\arraybackslash}p{(\linewidth - 6\tabcolsep) * \real{0.2381}}
  >{\raggedright\arraybackslash}p{(\linewidth - 6\tabcolsep) * \real{0.2619}}
  >{\raggedright\arraybackslash}p{(\linewidth - 6\tabcolsep) * \real{0.1905}}
  >{\raggedright\arraybackslash}p{(\linewidth - 6\tabcolsep) * \real{0.3095}}@{}}
\toprule\noalign{}
\begin{minipage}[b]{\linewidth}\raggedright
Strategy
\end{minipage} & \begin{minipage}[b]{\linewidth}\raggedright
Mechanism
\end{minipage} & \begin{minipage}[b]{\linewidth}\raggedright
Result
\end{minipage} & \begin{minipage}[b]{\linewidth}\raggedright
Failure mode
\end{minipage} \\
\midrule\noalign{}
\endhead
\bottomrule\noalign{}
\endlastfoot
Argmin (baseline) & \(g(s) = \text{argmin}_a N(s, a)\) & 10/10 & --- \\
Destination novelty & \(g(s) = \text{argmax}_a \text{novelty}(s'|s,a)\)
& 1/10 & Attracted to rare death states \\
Prediction error & \(g(s) = \text{argmax}_a \|s' - \hat{s}'\|\) & 0/10 &
Death is maximally unpredictable \\
Softmax temperature & \(P(a) \propto \exp(-N(s,a)/T)\) & 2/3 &
Stochasticity without benefit \\
Entropy-seeking & \(g(s) = \text{argmax}_a H(s'|s,a)\) & 0/3 & Noisy TV:
death = max entropy \\
UCB1 & \(g(s) = \text{argmax}_a (\bar{r}_a + c\sqrt{\ln t / N_a})\) &
2/3 & Degenerates to argmin (no reward) \\
Global novelty & \(g(s) = \text{argmax}_a (1/N_{\text{global}}(s'))\) &
6/10 & Same count, different seeds \\
\end{longtable}
}

Steps 477-482, 539-541. 6 strategies tested, all worse than or equal to
argmin.

\textbf{Why argmin is robust:} Define a transition as \emph{reducible}
if \(H(s' | s, a) \to 0\) with sufficient observations, and
\emph{irreducible} if \(H(s' | s, a) > \epsilon\) for all sample sizes
(e.g., death transitions, stochastic resets). Any strategy \(g\) that
selects actions based on model uncertainty (prediction error, entropy,
novelty) preferentially selects actions leading to irreducible
transitions, because these maximize the quality signal. Argmin
\(g(s) = \text{argmin}_a N(s, a)\) is immune: visit counts accumulate
equally on reducible and irreducible transitions. The cost is slower
exploration of structured regions; the benefit is avoiding traps.

\textbf{Post-ban resolution:} Compression progress (Schmidhuber 1991)
avoids noisy TV by rewarding learning RATE not error. Step 855: ACTION
COLLAPSE --- locks onto single action. Fix candidates untested.

\textbf{Proposition 15 (Perception-Action Decoupling).} L1 rate
determined by observation mapping \(\pi\), not action selection \(g\).
20+ interventions on \(g\) all fail; only \(\pi\) modifications improve
L1. → \url{propositions/15_perception_action_decoupling.md}

\textbf{Proposition 16 (Transition-Inconsistency Refinement).} Cells
with \(I(n) \geq 2\) conflate hidden states. Refining hash at these
cells: 17/20 L1 at 25s, FT09 5/5. →
\url{propositions/16_transition_inconsistency_refinement.md}

\textbf{L1/L2 Asymmetry.} L1 aliasing is bounded (83-313 cells),
solvable by refinement. L2 aliasing is unbounded, grows monotonically.
L2 requires prediction of unvisited states, not finer perception --- the
transition from \(\ell_\pi\) to \(\ell_F\).

\subsubsection{4.6 Constructive Characterization of the Feasible
Region}\label{constructive-characterization-of-the-feasible-region}

Combining all formalized constraints, we characterize the class of \(F\)
that satisfies R1-R6 + validated U-constraints simultaneously.

\textbf{The minimal frozen frame \(F\) decomposes as:}

\[F(s)(x) = \texttt{store}(s, \texttt{select}(s, \texttt{compare}(s, x)))\]

where: - \(\texttt{compare}(s, x)\): produces similarity/distance scores
between \(x\) and elements of \(s\) -
\(\texttt{select}(s, \text{scores})\): chooses an element of \(s\) (or
triggers creation of a new element) -
\(\texttt{store}(s, \text{selection})\): updates \(s\) with the result

\textbf{Constraints on each operation:}

The feasible region is the set of all \((compare, select, store)\)
triples satisfying all constraints. See MAP.md for the full constraint
table.

\subsubsection{4.7 Interpreter Entailment (Proposition
14)}\label{interpreter-entailment-proposition-14}

\textbf{The question (posed 2026-03-21):} Section 4.6 identifies the
minimal frozen frame as the interpreter: compare-select-store. When is
this interpreter \emph{entailed} by the system --- a logical consequence
of R1-R6 + task structure --- rather than a designer choice?

\textbf{Prior work:} Rosen's (M,R)-systems (1991) formalize
organizational closure, but Letelier et al.~(2006) showed the closure
cycle is closed by assumption, not derivation. Bauer (2016) proved
self-interpreters don't exist in total languages (System T).
Self-interpretation requires general recursion. Letelier et al.~(2023,
BioSystems) give partial solutions for restricted cases.

\textbf{Our formalization:} We distinguish two senses of ``entailment'':

\textbf{Proposition 14 (Interpreter Entailment).} Any self-referential
system satisfying R1-R6 + U3 has an irreducible top-level interpreter. →
\url{propositions/14_interpreter_entailment.md}

\textbf{Proposition 14b (CSE Uniqueness).} Compare-select-store is the
UNIQUE interpreter class at the coarsest decomposition for systems
satisfying R1-R6 + U3. → \url{propositions/14b_cse_uniqueness.md}

\textbf{Implication:} The frozen frame floor is the interpreter. R3
reduces to making everything ELSE (encoding, state, data)
self-modifiable --- the interpreter is fixed by mathematical necessity
(Proposition 12).

\subsubsection{4.8 R3 as Self-Directed Attention (Proposition
17)}\label{r3-as-self-directed-attention-proposition-17}

\textbf{The insight (compression, 2026-03-22):} R3 requires every aspect
of the system to be self-modified. CSE uniqueness (Proposition 14b) says
the interpreter structure is fixed. What remains to be self-modified?
The ENCODING --- the lens through which the interpreter sees. R3 is not
``modify the program.'' R3 is ``modify what the program attends to.''

\textbf{Prior work:} Active perception (Bajcsy, 1988) formalizes
perception as active selection, not passive reception. Friston (2010)
identifies attention as precision weighting --- the closest framework to
our ``self-directed attention,'' but adapts the entire generative
hierarchy rather than encoding alone. Predictive coding (Rao \& Ballard,
1999): centering (\(x - \mathbb{E}[x]\), U16) IS predictive coding with
the interpreter held fixed. Perceptual learning (Goldstone, 1998):
attentional weighting = D1, differentiation = D3. Efficient coding
(Barlow, 1961): we use \emph{transition} statistics rather than stimulus
statistics --- a more specific, action-driven signal.

\textbf{Proposition 17 (R3 as Self-Directed Attention).} For CSE
systems: R3 for comparison reduces to state-dependent encoding
\(\pi_s\). Not ``modify the program'' but ``modify the lens.'' →
\url{propositions/17_r3_self_directed_attention.md}

\textbf{Proposition 18 (Eigenform Reactivation).} Eigenform (inert for
action selection, Prop 13) reactivates for encoding adaptation via
transition statistics \(T(s)\). →
\url{propositions/18_eigenform_reactivation.md}

\textbf{Encoding Dimension Taxonomy.} Five dimensions of encoding
self-modification (D1-D5):

{\def\LTcaptype{none} % do not increment counter
\begin{longtable}[]{@{}
  >{\raggedright\arraybackslash}p{(\linewidth - 6\tabcolsep) * \real{0.1515}}
  >{\raggedright\arraybackslash}p{(\linewidth - 6\tabcolsep) * \real{0.3636}}
  >{\raggedright\arraybackslash}p{(\linewidth - 6\tabcolsep) * \real{0.2424}}
  >{\raggedright\arraybackslash}p{(\linewidth - 6\tabcolsep) * \real{0.2424}}@{}}
\toprule\noalign{}
\begin{minipage}[b]{\linewidth}\raggedright
Dim
\end{minipage} & \begin{minipage}[b]{\linewidth}\raggedright
What adapts
\end{minipage} & \begin{minipage}[b]{\linewidth}\raggedright
Signal
\end{minipage} & \begin{minipage}[b]{\linewidth}\raggedright
Status
\end{minipage} \\
\midrule\noalign{}
\endhead
\bottomrule\noalign{}
\endlastfoot
D1 & Channel weights & Per-channel prediction error & CONFIRMED (alpha,
Step 895) \\
D2 & Spatial resolution & Per-region transition inconsistency &
CONFIRMED (674, Step 690) \\
D3 & Hash K (partition granularity) & Aliased cell count & TESTED
(Recode, K confound) \\
D4 & Temporal depth & Temporal prediction improvement & UNTESTED \\
D5 & Centering rate & Running mean convergence & UNTESTED \\
\end{longtable}
}

\subsubsection{4.9 R3 Counterfactual Requirement (Proposition
19)}\label{r3-counterfactual-requirement-proposition-19}

\textbf{Prior work:} Gödel Machine (Schmidhuber, 2003) requires provable
improvement (intractable). Darwin Gödel Machine (Sakana AI, 2025)
relaxes to empirical improvement via evolutionary search. Off-policy
evaluation (Uehara et al., 2022) formalizes counterfactual comparison.
Our R3 counterfactual is empirical but R1-constrained (no external
objective).

\textbf{Our formalization:}

\textbf{Proposition 19 (R3 Counterfactual).} Graph state transfers
negatively: cold \textgreater{} warm (p\textless0.0001, Step 776). The
mechanism that makes navigation work IS the mechanism that prevents R3.
→ \url{propositions/19_r3_counterfactual.md}

\subsubsection{4.10 State Decomposition: Location vs Dynamics
(Proposition
20)}\label{state-decomposition-location-vs-dynamics-proposition-20}

\textbf{Prior work:} ExoMDP (Efroni et al., 2022) decomposes by
controllability; Denoised MDP (Wang et al., 2022) by reward-relevance;
Successor Features (Barreto et al., 2017) by transition structure vs
reward. None decompose along our location/dynamics axis relevant to
self-modification transfer. Bisimulation metrics (Zhang et al., 2021)
collapse both into one equivalence class.

\textbf{Proposition 20 (State Decomposition).} State decomposes into
location-dependent \(L(s)\) (visit counts --- always transfers
negatively) and dynamics-dependent \(D(s)\) (forward model --- can
transfer positively if dynamics generalize). →
\url{propositions/20_state_decomposition.md}

\textbf{Proposition 21 (Global-Local Gap).} \(D(s)\) captures global
dynamics; navigation needs local per-state selection. This gap is
structural --- no D-only mechanism consistently beats random for
navigation post-ban. → \url{propositions/21_global_local_gap.md}

\textbf{Key experimental evidence:} D(s) prediction transfer: 5/7 PASS
(Steps 778v5-855v3). Navigation transfer: 0 mechanisms beat random
consistently. The feasible region for prediction transfer is non-empty;
for navigation transfer it's empty.

\subsubsection{4.11 Architecture Triangle and R3 Structural Possibility
(Proposition
22)}\label{architecture-triangle-and-r3-structural-possibility-proposition-22}

\textbf{Proposition 22 (Architecture Triangle).} The 800+ experiments
across 12 architecture families cluster into three vertices defined by
what the substrate's state encodes:

\begin{enumerate}
\def\labelenumi{\arabic{enumi}.}
\item
  \textbf{Recognition vertex} (codebook family, Steps 1-416). State
  \(V = \{v_i\}\) encodes observation prototypes. Processing:
  \(i^* = \text{argmax}_i \cos(x, v_i)\),
  \(v_{i^*} \leftarrow v_{i^*} + \eta(x - v_{i^*})\). Action derives
  from nearest prototype. \textbf{R3 signal: none.} Cosine similarity is
  symmetric --- it measures proximity but cannot distinguish informative
  from uninformative observation dimensions.
\item
  \textbf{Tracking vertex} (graph family, Steps 417-777). State
  \(G: S \times A \to \mathbb{N}\) encodes transition counts.
  Processing: \(G(s, a) \leftarrow G(s, a) + 1\), action
  \(= \text{argmin}_a G(s, a)\). \textbf{R3 signal: successor set
  inconsistency} (Step 674 --- cells with \(|\text{succ}(n)| \geq 2\)
  trigger refinement). This IS an R3 signal (encoding modified by
  accumulated state). But the graph ban eliminates the data structure
  that generates it.
\end{enumerate}

\textbf{Proposition 22 (Architecture Triangle).} 800+ experiments
cluster into three vertices: recognition (codebook, banned), tracking
(graph, banned), dynamics (prediction, current frontier). Post-ban,
prediction error is the unique remaining signal for R3 encoding
self-modification. → \url{propositions/22_architecture_triangle.md}

\textbf{Corollary 22.1:} The true substrate lives at the dynamics vertex
--- the only vertex where circular causation (improve model → change
novelty → change exploration → improve model) creates (M,R)-system
closure.

\subsubsection{4.12 Game Taxonomy by Progress Structure (Proposition
23)}\label{game-taxonomy-by-progress-structure-proposition-23}

\textbf{Proposition 23 (Monotonic vs Sequential Progress).} Interactive
environments partition into two classes by progress structure:

\begin{enumerate}
\def\labelenumi{(\alph{enumi})}
\tightlist
\item
  \textbf{Observation-sufficient (monotonic) games.} Progress produces
  observable change: any action that changes the observation is
  (probabilistically) progressive. Navigation reduces to maximizing
  observation change. Change-tracking
  (\(\delta_a = \text{EMA}(\|\Delta x\|)\) per action) is sufficient.
  LS20 is observation-sufficient: moving the avatar changes the
  observation, and any movement direction that produces change leads
  toward unexplored territory.
\end{enumerate}

\textbf{Proposition 23 (Game Taxonomy by Progress Structure).} Games
split into monotonic (LS20: progress visible in observations) and
sequential (FT09: 7-step ordered puzzle, progress invisible until
completion). Global-EMA selectors work on monotonic games only. →
\url{propositions/23_game_taxonomy.md}

\textbf{Proposition 23b (Combinatorial Barrier).} Global-EMA selectors
are position-blind:
\(P(\text{correct 7-step sequence}) \leq (1/68)^6 \approx 10^{-11}\)/attempt.
Structurally unsolvable within 10K budget. →
\url{propositions/23b_combinatorial_barrier.md}

\subsubsection{4.13 Active Inference as the Action Selection Framework
(Proposition
24)}\label{active-inference-as-the-action-selection-framework-proposition-24}

\textbf{Prior work:} Active inference (Friston, 2009) replaces reward
with free energy minimization. Under R1, EFE reduces to pure epistemic
value (Sajid et al., 2021; Da Costa et al., 2020).

\textbf{Proposition 24:} Epistemic value
\(G_t(a) = \|\hat{x}^{(a)} - x_t\| / \text{conf}_a\) dissolves Prop 23
in theory. FAILED (Steps 934-936): \(W\) errors overwhelm signal. →
\url{propositions/24_active_inference.md}

\textbf{Proposition 25:} Adaptive EMA decay. FAILED (Step 937): signal
degradation without exact 916 formula. →
\url{propositions/25_adaptive_ema.md}

\subsubsection{4.14 Action Selection Closure (Steps
938-938e)}\label{action-selection-closure-steps-938-938e}

\textbf{Proposition 26 (Novelty-Reactive Policy).} REJECTED --- gate 5
(per-observation conditioning). Per-observation policy table stores
action per observation hash = per-observation conditioning, same
violation as Step 931 (PB20). →
\url{propositions/26_novelty_reactive_policy.md}

\textbf{Steps 938b-938e:} Four remaining mechanism classes tested
(reactive-global, trajectory-conditioned ×2, trajectory-additive). All
killed. \(h\) contains real position-dependent variance but mapping
\(h \to\) actions requires learned weights (R1 violation), W-based
projection (killed), or per-observation memory (gate 5). No
gate-compliant mapping exists.

\textbf{Conclusion:} Under gates 3-5 + bans, 800b is the unique working
action selector. The degree of freedom is closed.

\subsubsection{4.15 Growing Feature Space (Proposition 27, Step
939)}\label{growing-feature-space-proposition-27-step-939}

Online PCA grows encoding from observation statistics. →
\url{propositions/27_growing_feature_space.md}

\textbf{KILLED (Step 939).} PCA features discovered (10 in LS20, 16 by
FT09) but zero-initialized \(W_{pred}\) rows → huge prediction errors →
\(\alpha\) concentrates on new dims → navigation signal drowned out.

\subsubsection{4.16 Alpha Concentration Is Load-Bearing (Proposition 28
---
FALSIFIED)}\label{alpha-concentration-is-load-bearing-proposition-28-falsified}

\textbf{Proposition 28: FALSIFIED (Step 944).} Threshold reset at ALL
values degrades LS20. Seeds with HIGH alpha\_conc navigate WELL.
Concentration on informative dimensions IS the mechanism --- resetting
destroys learned attention weighting.

\textbf{R4 clarification:} Discriminative capacity (Ashby) is about
WHERE alpha concentrates, not WHETHER. alpha\_conc=50 on informative
dims = correct; on uninformative dims (GFS case, zero-init W\_pred) =
degenerate. The metric alone can't distinguish. Encoding expansion
requires warm-up before alpha integration.

\subsubsection{4.17 Temporal Credit Impossibility (Theorem 4,
Proposition
31)}\label{temporal-credit-impossibility-theorem-4-proposition-31}

\textbf{Prior work:} Temporal credit assignment (survey: Ferret et al.,
2023, arXiv:2312.01072) is defined as connecting actions to delayed
consequences. All standard mechanisms require some form of
state-conditioned memory: eligibility traces (Sutton, 1988) maintain
per-(state,action) decay traces; empowerment (Klyubin et al., 2005;
Salge et al., 2014) computes
\(\mathcal{E}(s) = I(A_{1:n}; S_{t+n} | S_t = s)\) requiring per-state
action sampling; PAC-MDP algorithms (Strehl \& Littman, 2008) require
visit counts. Phasor Agents (Trappe, 2026) use oscillatory phase
alignment for local unsupervised credit --- the closest mechanism to
reward-free temporal credit, but require per-edge eligibility traces.

\textbf{Our formalization:} Let
\(\bar{\Delta}(a) = \text{EMA}_t(\|f_\theta(s_t) - s_{t+1}\|)\) be the
global running mean of prediction error for action \(a\), averaged over
ALL states where \(a\) was executed.

\textbf{Theorem 4 (Temporal Credit Impossibility under Graph Ban).} For
environments where the optimal action depends on state
(\(\exists s \neq s': a^*(s) \neq a^*(s')\)), the signal-to-noise ratio
for identifying state-dependent optimal actions via \(\bar{\Delta}(a)\)
satisfies:

\[\text{SNR}(a, s^*) \leq \frac{1}{|N_a|}\]

where \(|N_a|\) is the number of distinct states at which action \(a\)
has been executed. Under uniform exploration, \(|N_a| \propto T/|A|\)
grows with time \(\Rightarrow\) SNR \(\to 0\).

\emph{Proof:}
\(\bar{\Delta}(a) = \frac{1}{|N_a|}\sum_{s \in N_a} \Delta(a|s)\). The
``correct'' state \(s^*\) contributes signal \(\delta^*\); all others
contribute mean \(\mu \pm \sigma\). The signal \(\delta^* - \mu\) is
diluted by \(1/|N_a|\) in the average. For LS20 (\(|A|=4\)): each
action's effect is state-independent (directional movement) →
\(\delta^* \approx \mu\) → SNR ≈ 1. For FT09 (\(|A|=68\)): action
effects are state-dependent → \(|N_a| \sim 150\) at 10K → SNR
\(\sim 0.007\). \(\square\)

\textbf{Corollary 4.1:} Eligibility traces require per-(state,action)
data (graph-banned). Empowerment requires per-state action sampling
(graph-banned). Both standard solutions to temporal credit are
prohibited. \textbf{No known mechanism provides temporal credit for
ordered sequences under the graph ban without external reward.}

\textbf{Relationship to prior work:} The SNR dilution of global
statistics is related to state aliasing in POMDPs (McCallum, 1996) and
the PAC-MDP impossibility without visit counts (Strehl \& Littman,
2008). Our specific contribution: (1) proving the impossibility bound
for the global running mean mechanism, (2) connecting it to the graph
ban as the structural cause, and (3) the 43-experiment empirical
confirmation across 5+ mechanism families (Steps 948-990).

\textbf{Degrees of freedom:} The theorem identifies what MUST change for
FT09/VC33: either (a) a non-running-mean mechanism for temporal credit
without per-state data, (b) lifting the graph ban (parametric models ---
but Step 972 shows these compute the same function), or (c) a mechanism
class outside prediction-error exploration entirely.

\subsubsection{4.18 Component Extraction Validity (Theorem
5)}\label{component-extraction-validity-theorem-5}

\textbf{Prior work:} SURD (Martínez-Sánchez, Arranz, Lozano-Durán,
Nature Communications 2024) decomposes causality into synergistic,
unique, and redundant components. Ablation methodology (Meyes et al.,
2019) establishes that removal ablation (component omitted during
training and testing) is the valid paradigm; inaccurate ablation
(corrupted at test only) introduces confounds. Our extraction protocol
is a form of removal ablation applied to banned architecture families.

\textbf{Our formalization:} Let family
\(\mathcal{F} = \{c_1, \ldots, c_n\}\) be a system of components, and
let \(P\) be the property identified by the killing finding \(K\):
\(K(\mathcal{F}) = \text{FAIL}\) because
\(P(\mathcal{F}) = \text{true}\).

Define the causal contribution types of component \(c_i\) to property
\(P\): - \textbf{Unique:} \(P(\{c_i\} \cup S') = \text{true}\) for any
substrate \(S'\) not containing other \(\mathcal{F}\) components.
\(c_i\) alone carries \(P\). - \textbf{Synergistic:}
\(P(\{c_i\} \cup S') = \text{false}\) when
\(\{c_j, c_k, \ldots\} \not\subseteq S'\). \(P\) requires the
combination, not \(c_i\) alone. - \textbf{Redundant:}
\(P(\mathcal{F} \setminus \{c_i\}) = \text{true}\). \(P\) persists even
without \(c_i\) (other components also produce \(P\)).

\textbf{Theorem 5 (Extraction Protocol Validity).} The extraction
protocol --- testing \(S' \cup \{c_i\}\) against killing finding \(K\)
--- correctly identifies whether \(c_i\) carries property \(P\) if and
only if \(c_i\)'s causal contribution to \(P\) is unique or redundant.
For synergistic contributions, the protocol produces false negatives:
\(c_i\) appears safe when tested alone, but \(P\) re-emerges if
synergistic partners are later added.

\emph{Proof sketch:} If \(c_i\)'s contribution is unique,
\(S' \cup \{c_i\}\) has \(P\) regardless of \(S'\)'s composition → test
detects \(P\) → correctly rejects \(c_i\). If redundant, same logic
applies. If synergistic with \(\{c_j, c_k\}\) and
\(\{c_j, c_k\} \not\subseteq S'\), then
\(P(S' \cup \{c_i\}) = \text{false}\) → test clears \(c_i\) → but
\(P(S' \cup \{c_i, c_j, c_k\}) = \text{true}\). \(\square\)

\textbf{Corollary 5.1 (Safe extraction criterion).} Component \(c_i\)
can be safely extracted from banned family \(\mathcal{F}\) if: (a) the
killing finding's property \(P\) is localized to a different component
\(c_j \neq c_i\) (unique to \(c_j\)), or (b) \(c_i\) serves a function
orthogonal to \(P\) (e.g., \(c_i\) is an encoding component and \(P\) is
about action selection).

\textbf{Application to our bans:}

{\def\LTcaptype{none} % do not increment counter
\begin{longtable}[]{@{}
  >{\raggedright\arraybackslash}p{(\linewidth - 8\tabcolsep) * \real{0.0725}}
  >{\raggedright\arraybackslash}p{(\linewidth - 8\tabcolsep) * \real{0.3043}}
  >{\raggedright\arraybackslash}p{(\linewidth - 8\tabcolsep) * \real{0.1594}}
  >{\raggedright\arraybackslash}p{(\linewidth - 8\tabcolsep) * \real{0.1884}}
  >{\raggedright\arraybackslash}p{(\linewidth - 8\tabcolsep) * \real{0.2754}}@{}}
\toprule\noalign{}
\begin{minipage}[b]{\linewidth}\raggedright
Ban
\end{minipage} & \begin{minipage}[b]{\linewidth}\raggedright
Killing property \(P\)
\end{minipage} & \begin{minipage}[b]{\linewidth}\raggedright
Component
\end{minipage} & \begin{minipage}[b]{\linewidth}\raggedright
Causal type
\end{minipage} & \begin{minipage}[b]{\linewidth}\raggedright
Extraction valid?
\end{minipage} \\
\midrule\noalign{}
\endhead
\bottomrule\noalign{}
\endlastfoot
Codebook & LVQ characterization & Attract dynamics & Synergistic (with
cosine) & \textbf{Risky} --- LVQ needs all three \\
Codebook & LVQ characterization & Novelty spawning & Unique to growth &
\textbf{Valid} --- spawning ≠ LVQ \\
Codebook & LVQ characterization & Centered encoding & Unique to
preprocessing & \textbf{Valid} --- already U16 \\
Graph & Negative transfer & Per-state-action storage & Unique to storage
& \textbf{Valid} --- this IS \(P\) \\
Graph & Negative transfer & Transition detection & Unique to detection &
\textbf{Valid} --- detection ≠ storage \\
Graph & Negative transfer & CC zone discovery & Unique to encoding &
\textbf{Valid} --- discovery ≠ per-state data \\
\end{longtable}
}

\textbf{Relationship to prior work:} SURD provides the
information-theoretic framework for continuous time series; we apply to
discrete component systems. Ablation methodology (Meyes et al., 2019)
addresses removal in neural networks; we extend to cross-family
transplantation. The extraction protocol formalization is novel.

\textbf{Degrees of freedom:} The theorem doesn't tell us which
extraction tests will SUCCEED (produce useful new substrates) --- only
which are VALID (correctly test what they claim to test). 33 components
cataloged (COMPONENT\_CATALOG.md); 5 extraction experiments designed
with valid causal isolation (experiments/extraction\_specs.md).

\subsection{5. Experimental Evidence}\label{experimental-evidence}

\subsubsection{5.1 Navigation (943+
experiments)}\label{navigation-943-experiments}

All 3 ARC-AGI-3 games solved at Level 1. Unifying mechanism: graph +
edge-count argmin + correct action/encoding decomposition.

{\def\LTcaptype{none} % do not increment counter
\begin{longtable}[]{@{}
  >{\raggedright\arraybackslash}p{(\linewidth - 4\tabcolsep) * \real{0.3333}}
  >{\raggedright\arraybackslash}p{(\linewidth - 4\tabcolsep) * \real{0.3333}}
  >{\raggedright\arraybackslash}p{(\linewidth - 4\tabcolsep) * \real{0.3333}}@{}}
\toprule\noalign{}
\begin{minipage}[b]{\linewidth}\raggedright
Game
\end{minipage} & \begin{minipage}[b]{\linewidth}\raggedright
Best result
\end{minipage} & \begin{minipage}[b]{\linewidth}\raggedright
Key mechanism
\end{minipage} \\
\midrule\noalign{}
\endhead
\bottomrule\noalign{}
\endlastfoot
LS20 & 20/20 at 25s (674+running-mean) & LSH dual-hash, 4 actions \\
FT09 & 3/3 at 50K (k-means, 69 actions) & Grid-click expansion (Step
503) \\
VC33 & 3/3 at 30K (k-means, 3 actions) & Zone discovery (Step 505) \\
\end{longtable}
}

\textbf{POMDP reframing (Steps 652-690).} L1 is a hidden-state
conjunction problem --- the agent visits the exit cell avg 152 times
before L1 triggers (Prop 15). Step 674 (transition-triggered dual-hash):
cells with inconsistent transitions get finer hash → 17/20 at 25s, 20/20
at 120K. This is \(\ell_\pi\): encoding refined from transition
statistics, R1-compliantly.

\textbf{Encoding requirements:} centered\_enc load-bearing across 2
families. avgpool16 required for LS20 at k=16. See CONSTRAINTS.md for
full evidence.

\subsubsection{5.2 Level 2 Failure as Feasibility
Violation}\label{level-2-failure-as-feasibility-violation}

L2 failure is a feasibility violation, not a performance gap. The
substrate's state graph at 500K steps: 942 live cells, 364-node active
set (29\%), 134 unexploited frontier edges. Extended budget produces no
improvement --- active set plateaus at \textasciitilde5000 (Step 555).
Root cause: energy mechanic requires navigating TO energy sources before
dying (Steps 556-557). Rudakov et al.~(2025) solved L2 using
connected-component segmentation with visual salience prioritization.

\subsubsection{5.3 Architecture Family
Summary}\label{architecture-family-summary}

{\def\LTcaptype{none} % do not increment counter
\begin{longtable}[]{@{}
  >{\raggedright\arraybackslash}p{(\linewidth - 8\tabcolsep) * \real{0.1739}}
  >{\raggedright\arraybackslash}p{(\linewidth - 8\tabcolsep) * \real{0.2609}}
  >{\raggedright\arraybackslash}p{(\linewidth - 8\tabcolsep) * \real{0.1087}}
  >{\raggedright\arraybackslash}p{(\linewidth - 8\tabcolsep) * \real{0.2826}}
  >{\raggedright\arraybackslash}p{(\linewidth - 8\tabcolsep) * \real{0.1739}}@{}}
\toprule\noalign{}
\begin{minipage}[b]{\linewidth}\raggedright
Family
\end{minipage} & \begin{minipage}[b]{\linewidth}\raggedright
Experiments
\end{minipage} & \begin{minipage}[b]{\linewidth}\raggedright
L1
\end{minipage} & \begin{minipage}[b]{\linewidth}\raggedright
Key finding
\end{minipage} & \begin{minipage}[b]{\linewidth}\raggedright
Status
\end{minipage} \\
\midrule\noalign{}
\endhead
\bottomrule\noalign{}
\endlastfoot
Codebook/LVQ & 416 & Yes (26K) & IS LVQ (Kohonen 1988). Fully
characterized. & \textbf{BANNED} (Step 416) \\
LSH & 100+ & Yes (674: 20/20) & Dual-hash + transition refinement. Best
L1. & ACTIVE (frozen bootloader) \\
Graph & 50+ & Yes (argmin) & Architecture-invariant argmin. Three
mapping props. & \textbf{BANNED} (Step 777) \\
Reservoir & 8 & No & Rank-1 collapse (U7). & KILLED \\
Hebbian & 4 & No & W diverges. Delta rule stabilizes but no nav. &
KILLED \\
Recode & 33 & Yes (5/5) & Self-refinement achieves \(\ell_\pi\). K
confound. & SUSPENDED \\
SplitTree & 3 & No & Binary splits too coarse for continuous obs. &
KILLED \\
Absorb & 2 & No & Merge rule too aggressive. & KILLED \\
ConnComp & 5 & No & Segmentation noise. & KILLED \\
Bloom & 2 & No & False positives compound. & KILLED \\
CA & 3 & No & Rule space too large for online search. & KILLED \\
800b-variants & 160+ & Partial (LS20 only) & Position-blind.
Classification-blind. & \textbf{KILLED} (Step 937) \\
\end{longtable}
}

\subsubsection{5.4 The Chain Benchmark}\label{the-chain-benchmark}

The chain benchmark (Split-CIFAR-100 → LS20 → FT09 → VC33 →
Split-CIFAR-100) tests the substrate on the FULL trajectory:
classification + 3 navigation games + classification again. One
configuration, no resets between tasks.

\textbf{Key findings (Steps 506-546):} - Negative transfer: CIFAR
pre-training degrades LS20 (Step 506-508, 14\% L1 regression) -
Threshold tension: optimal game thresholds differ by 2.5x (Step 509-513)
- Algorithm invariance: 4 representations (edge dict, W matrix, tensor,
n-gram) produce identical L1 rates (Steps 521-525) - Recode
self-refinement: LSH k=16 + self-refinement = 5/5 L1 (Step 542). K
confound invalidates some results (Step 589).

\textbf{Chain kill criterion (Jun, 2026-03-23):} Any mechanism that
improves one game at the cost of another = per-game tuning = KILL. Only
mechanisms neutral/positive on ALL games survive.

\subsubsection{5.5 Post-Ban Results (Steps 778-1007, 230+
experiments)}\label{post-ban-results-steps-778-1007-230-experiments}

Both working mechanisms banned (codebook Step 416, graph Step 777). Key
results: - D(s) prediction transfer: 5/7 PASS --- first positive R3\_cf
(+73\%, Step 780v5) - 800b (per-action change EMA): best post-ban
mechanism. LS20=268.0/seed cold (895h), 290.7/seed (916 with recurrent
h). FT09/VC33=0 (Theorem 4). - Alpha prediction-error attention: R3
encoding confirmed. FT09 dims {[}60,51,52{]} discovered autonomously. -
800b action selection: FROZEN. 27 modifications all kill LS20
(kills/800b-variants\_step937.md). - Extraction protocol: 33 components
cataloged, 5 extraction experiments designed (Section 4.18).

\subsubsection{5.7 Baseline Comparison (Steps
916-920)}\label{baseline-comparison-steps-916-920}

Published baselines reproduced in our framework on LS20 (25K steps,
\(n_{\text{eff}}=10\), substrate\_seed=seed):

{\def\LTcaptype{none} % do not increment counter
\begin{longtable}[]{@{}
  >{\raggedright\arraybackslash}p{(\linewidth - 8\tabcolsep) * \real{0.1702}}
  >{\raggedright\arraybackslash}p{(\linewidth - 8\tabcolsep) * \real{0.1915}}
  >{\raggedright\arraybackslash}p{(\linewidth - 8\tabcolsep) * \real{0.1064}}
  >{\raggedright\arraybackslash}p{(\linewidth - 8\tabcolsep) * \real{0.2553}}
  >{\raggedright\arraybackslash}p{(\linewidth - 8\tabcolsep) * \real{0.2766}}@{}}
\toprule\noalign{}
\begin{minipage}[b]{\linewidth}\raggedright
Method
\end{minipage} & \begin{minipage}[b]{\linewidth}\raggedright
L1/seed
\end{minipage} & \begin{minipage}[b]{\linewidth}\raggedright
std
\end{minipage} & \begin{minipage}[b]{\linewidth}\raggedright
zero seeds
\end{minipage} & \begin{minipage}[b]{\linewidth}\raggedright
Signal type
\end{minipage} \\
\midrule\noalign{}
\endhead
\bottomrule\noalign{}
\endlastfoot
\textbf{916 Recurrent \(h\)} & \textbf{290.7} & \textbf{70.1} &
\textbf{0/10} & pred error + trajectory \\
\textbf{895h cold} & \textbf{268.0} & \textbf{75.2} & \textbf{0/10} &
pred error (\(\alpha\)-weighted \(\Delta\)) \\
868d raw L2 & 203.9 & 105.8 & 1/10 & raw \(\Delta\) \\
920 Graph+argmin & 129.9 & 124.0 & 4/10 & visit count \\
918 RND (Burda 2018) & 112.3 & 126.0 & 4/10 & distillation error \\
919 Count-based (Bellemare 2016) & 109.0 & 125.3 & 5/10 & obs
frequency \\
917 ICM (Pathak 2017) & 0.0 & 0.0 & 10/10 & forward pred error \\
\end{longtable}
}

All methods L1=0 on FT09 (68 actions). Even graph+argmin at 6 correct
actions + avgpool16 = L1=0 (Step 920b). Generic exploration insufficient
--- FT09 6/6 required per-game prescribed solution (Step 1012, all
constraints lifted). Bottleneck is solution architecture (prescribed vs
discovered), not encoding or action count.

\textbf{Finding:} Our mechanism outperforms ALL baselines by
\(2\text{-}2.5\times\) on LS20 with 0/10 zero-seeds. Graph ban cost is
NEGATIVE: 895h (268.0, no graph) \(>\) 920 (129.9, with graph). ICM
worst (signal dies as \(W\) learns).

\subsubsection{5.8 Chain and PRISM
Results}\label{chain-and-prism-results}

\textbf{PRISM baseline (Step 1006):} CIFAR=100\%, LS20=100\%, FT09=0\%,
VC33=0\%, chain\_score=3/5. CIFAR inflates alpha\_conc (PB16, −11\%
chain LS20). Delta inversion (Remark 5.2): 800b anti-selects on
reset-heavy games --- structural, not fixable. \textbf{Research skew:}
230 post-ban experiments exclusively LS20. FT09/VC33 = 0 post-ban
experiments. Extraction sprint (Steps 1008+) corrects this.

\subsubsection{5.9 Unconstrained Diagnostic (Steps
1012-1017)}\label{unconstrained-diagnostic-steps-1012-1017}

\textbf{Prescribed:} FT09 6/6 (Step 1012), VC33 7/7 (Step 1013).
\textbf{Generic with ALL bans lifted (Step 1017):} 674 graph + CC zones
→ FT09=0\%, VC33=0\%. The 13-level gap is PRESCRIPTION, not constraints.
Temporal credit for multi-step sequences within tight budgets --- not
state coverage --- is the bottleneck. No mechanism in 1000+ experiments
addresses this.

\subsection{6. Degrees of Freedom}\label{degrees-of-freedom}

The formalization identifies what the constraints REQUIRE but also what
they leave UNDETERMINED. These degrees of freedom define the experiment
space for the next phase.

{\def\LTcaptype{none} % do not increment counter
\begin{longtable}[]{@{}
  >{\raggedright\arraybackslash}p{(\linewidth - 8\tabcolsep) * \real{0.2000}}
  >{\raggedright\arraybackslash}p{(\linewidth - 8\tabcolsep) * \real{0.2000}}
  >{\raggedright\arraybackslash}p{(\linewidth - 8\tabcolsep) * \real{0.2000}}
  >{\raggedright\arraybackslash}p{(\linewidth - 8\tabcolsep) * \real{0.2000}}
  >{\raggedright\arraybackslash}p{(\linewidth - 8\tabcolsep) * \real{0.2000}}@{}}
\toprule\noalign{}
\begin{minipage}[b]{\linewidth}\raggedright
\#
\end{minipage} & \begin{minipage}[b]{\linewidth}\raggedright
Degree of Freedom
\end{minipage} & \begin{minipage}[b]{\linewidth}\raggedright
What's constrained
\end{minipage} & \begin{minipage}[b]{\linewidth}\raggedright
What's free
\end{minipage} & \begin{minipage}[b]{\linewidth}\raggedright
Source
\end{minipage} \\
\midrule\noalign{}
\endhead
\bottomrule\noalign{}
\endlastfoot
1 & Growth function \(\phi\) & Must be monotonic, unbounded (U17) & What
grows --- nodes, edges, meta-state, or something else & Sec 3.2 \\
2 & Oscillation period & System cannot converge (Thm 1) & How quickly
growth disrupts the dominant mode & Sec 4.1 \\
3 & Meta-rule \(F\) & Must be non-constant (R3), time-invariant (U1),
minimal (R6) & The specific form of compare-select-store & Sec 3.2 \\
4 & Growth-rule coupling & Growth changes the update rule (R3) & Local
vs global: does a new node change nearby rules or all rules? & Sec
3.3 \\
5 & Metric \(d_X\) on observations & Must make \(\pi\) Lipschitz (U20);
can refine but not rearrange (T3) & Which metric: L2, cosine, learned,
adaptive & Sec 3.3 \\
6 & Partition resolution & Nodes must be well-separated (Prop 2) & How
many nodes: determined by spawn threshold vs observation distribution &
Sec 3.3 \\
7 & Growth mechanism & Structure only grows (U3), without bound (U17) &
Spawn-on-novelty vs implicit edge accumulation vs other & Sec 3.3 \\
8 & Self-observation target & Must extract irredundant structure from
\(s\) (Thm 2) & What to extract: temporal patterns, graph properties,
meta-state & Sec 4.3 \\
9 & Self-observation → rule modification & New structure must change
\(f\) (R3) & How: does detecting stuckness change exploration? Does
cycle detection trigger avoidance? & Sec 4.3 \\
10 & Fixed-point structure & Conjectured to exist (Sec 4.4) & What
invariant is preserved under self-observation & Sec 4.4 \\
11 & Exploration/exploitation resolution & \(g_{nav} \neq g_{class}\)
(U11), no mode switch (U1) & State-dependent behavior: what state
feature triggers the switch? & Sec 3.4 \\
12 & Self-modification level & Must be \(\geq \ell_\pi\) for R3
compliance (Sec 3.2) & How deep: \(\ell_\pi\) (Recode) vs \(\ell_F\)
(full self-reference). Which component of \(F\) adapts, and by what
trigger? & Sec 3.2 \\
13 & Encoding channel selection & Must preserve transition signal
(R1-compliant) & Which channels: greyscale, RGB, subset? Weights from
per-channel transition discrimination & Sec 4.8 (D1) \\
14 & Encoding spatial resolution & Must be Lipschitz (U20); finer is not
always better (Step 597: K=16/10K, random beats argmin) & Pool size per
region from local transition inconsistency & Sec 4.8 (D2) \\
15 & Encoding temporal depth & Single frame or multi-frame & Frame stack
depth from temporal autocorrelation at each cell & Sec 4.8 (D4) \\
16 & Encoding-statistics coupling rate & Must be non-zero (R3) and
finite (stability) & How quickly encoding parameters respond to
transition statistics & Sec 4.8 (T9, DoF 12) \\
\end{longtable}
}

\textbf{Central experimental question:} R3 = self-directed attention
(Props 14b, 17, 18). Search space: encoding dimensions (D1-D5, §4.8)
made state-dependent via transition statistics. Each independently
testable on the chain benchmark.

\subsection{7. Discussion}\label{discussion}

\subsubsection{7.1 Pairwise Consistency
Audit}\label{pairwise-consistency-audit}

9 tensions checked. 6 resolved (T2: R3≡attention, T4: irredundant
growth, T5: state-derived components, T6: one rule different inputs, T7:
count-based self-observation, T8: per-domain centering). 1 constrained
(T3: metric refines not rearranges). 2 open (T1: U7 reformulation, T9:
encoding-statistics feedback). No contradictions found.

\subsubsection{7.2 What is proven}\label{what-is-proven}

\begin{enumerate}
\def\labelenumi{\arabic{enumi}.}
\tightlist
\item
  \textbf{Theorem 1} (No Fixed Point): R3 + U7 + U17 + U22 → no
  satisfying system has a fixed point.
\item
  \textbf{Theorem 2} (Self-Observation): In finite environments,
  irredundant growth requires processing internal state.
\item
  \textbf{Theorem 3} (System Boundary): No system can modify its own
  ground truth test.
\item
  \textbf{Propositions 14/14b:} CSE is the unique irreducible
  interpreter for R1-R6 + U3 systems.
\item
  \textbf{Proposition 15:} L1 rate is determined by \(\pi\), not \(g\).
  20+ interventions confirm.
\item
  \textbf{Proposition 16:} Transition inconsistency detects aliased
  cells. 674: 17/20 L1 at 25s.
\item
  \textbf{Proposition 17:} R3 = self-directed attention (encoding, not
  interpreter).
\item
  \textbf{Proposition 19:} Graph state transfers negatively (cold
  \textgreater{} warm, p\textless0.0001).
\item
  \textbf{Proposition 20:} State = L(s) + D(s). L always transfers
  negatively. D can transfer positively.
\item
  \textbf{Alpha encoding self-modification:} R3\_dynamic=1.0 across
  games (Steps 895-895h). CONFIRMED.
\item
  \textbf{Theorem 4} (Temporal Credit Impossibility): Global running
  mean SNR → 0 for state-dependent actions under graph ban. 43
  experiments confirm (Steps 948-990).
\end{enumerate}

\subsubsection{7.3 What is conjectured}\label{what-is-conjectured}

\begin{enumerate}
\def\labelenumi{\arabic{enumi}.}
\setcounter{enumi}{11}
\tightlist
\item
  The minimal self-observing substrate is an eigenform of \(F\): \(F\)
  applied to its own state reproduces its dynamics (Section 4.4).
\item
  The oscillation between U7 (convergence) and U17 (growth) is a limit
  cycle, not chaos.
\item
  Selection pressure on a population of encodings (GRN architecture) may
  bridge the ℓ₁→ℓ\_F gap without optimization. Step 607 tested: KILLED
  (Cov(w,z)=0). Framework sound but requires behaviorally diverse
  encodings.
\item
  \textbf{Search reduction (Propositions 14b + 17 + 18 combined).} The
  path to R3 is not building a new interpreter --- it's building
  self-directed attention into the existing CSE interpreter. The search
  space reduces from ``all self-modifying dynamical systems'' to ``5
  encoding dimensions (D1-D5) that can be made state-dependent via
  transition statistics.'' Each dimension is independently testable in 5
  minutes. Conditional on CSE uniqueness and hierarchy collapse holding.
\end{enumerate}

\subsubsection{7.4 What is open}\label{what-is-open}

\begin{enumerate}
\def\labelenumi{\arabic{enumi}.}
\setcounter{enumi}{14}
\tightlist
\item
  \textbf{Full feasible region occupancy.} Theorem 3 shows not provably
  empty; witness or tighter impossibility needed.
\item
  \textbf{Self-observation mechanism.} Must satisfy R6 while avoiding
  noisy TV traps. Compression progress fails (Step 855, action
  collapse). Prop 17: self-directed attention via \(\Delta E_t\) is
  R6-compliant.
\item
  \textbf{R1-compliant classification.} 0 substrates above chance
  without external labels (best: LSH k=16, 36.2\% self-labels, Step
  573).
\item
  \textbf{Post-ban action selection.} 800b is the unique working
  selector. FT09 = 0 formally characterized by Theorem 4 (temporal
  credit impossibility). → kills/800b-variants\_step937.md
\item
  \textbf{Post-ban constraint map.} Many U-constraints derived within
  graph framework. Survival post-ban is open.
\item
  \textbf{Component extraction.} Ban extraction protocol (2026-03-24)
  permits isolated mechanism testing. 33 components cataloged across 14
  families (COMPONENT\_CATALOG.md). Priority: FT09/VC33 action-space
  discovery components (CC zone discovery, transition detection, sparse
  gating).
\item
  \textbf{LS20 research skew.} 230 post-ban experiments focused
  exclusively on LS20 (solvable). FT09/VC33 abandoned since graph ban
  (Step 777). The search has TWO distinct problems: (a) exploration
  (LS20, well-characterized by 500+ experiments), (b) action-space
  discovery (FT09/VC33, 0 post-ban experiments). The extraction protocol
  enables Problem (b) experiments for the first time.
\end{enumerate}

\subsubsection{7.5 The Level 2 Problem}\label{the-level-2-problem}

L2 requires: level-aware background modeling, robust object detection,
hidden-state coverage (POMDP with \(|S| \leq 96\)). Systematic
falsification cascade (Steps 571-572j, 10 iterations): each fix exposed
the next blocker. Final: dead reckoning + state estimation + sequencing
= L2=5/5 at 4804 steps (Step 572j). The cascade reveals 12 orthogonal
prescribed layers constituting the L2 frozen frame.

\subsubsection{7.6 Honest assessment}\label{honest-assessment}

The feasible region for L1 navigation is occupied --- graph + argmin +
correct encoding solves all three games. But argmin's advantage over
random is speed, not exclusive access (13/20 vs 10/20, p=0.26). L2 via
prescribed 12-component pipeline (Step 572j, 5/5 at 4804 steps).
Pipeline is R1-compliant in detection but game-specific in state
estimation.

\textbf{R3 gap:} 12 design choices the substrate cannot self-discover.
Four general techniques, eight game-specific. A substrate satisfying R3
would discover all 12 from interaction alone.

\textbf{Proposition 4 (R3-R5 Tension).} R3 demands change; R5 demands
stability. Resolution: the frozen frame is the empirical ground truth
test, minimal by R6. → \url{propositions/04_r3_r5_tension.md}

\textbf{Theorem 3 (System Boundary).} No system can modify its own
ground truth test --- that changes the metric, not the system. →
\url{propositions/theorem1_no_fixed_point.md}

\textbf{Proposition 6 (Self-Modification Level).} Higher modification
levels trade speed for reachability. \(\ell_0\) (fixed rules) is fastest
but ceiling-limited. \(\ell_F\) (rule modification) is slowest but can
escape any fixed ceiling. →
\url{propositions/06_self_modification_level.md}

\textbf{Proposition 9 (R1 Tax).} R1-compliant systems cannot access
supervised signals. Classification without external labels = chance.
Navigation (self-generated actions) IS R1-compliant. →
\url{propositions/09_r1_tax.md}

\textbf{Proposition 10 (Feature-Ground Truth Coupling).} Features
learned from observations couple to the ground truth's measurement
scale. The system cannot escape this coupling without modifying the
ground truth (blocked by R5). →
\url{propositions/10_feature_ground_truth_coupling.md}

\textbf{Proposition 11 (Frozen Frame Capability Coupling).} The frozen
frame IS the navigation capability. Modifying it destroys navigation. R3
and navigation are in tension --- relaxing one strengthens the other. →
\url{propositions/11_frozen_frame_capability_coupling.md}

\textbf{Proposition 12 (Interpreter Bound).} Every self-referential
system has an irreducible top-level interpreter. The frozen frame floor
\textgreater{} 0. But \(\ell_F\) achievable if state encodes operations
the interpreter executes. → \url{propositions/12_interpreter_bound.md}

\textbf{Post-ban addendum (Steps 778-938e):} 800b achieves 2x random on
LS20. FT09 = 0 for all mechanisms. CIFAR = chance. R3 encoding
self-modification CONFIRMED (alpha, Step 895). Action selection CLOSED
after 938 series: per-observation (gate 5), reactive-global (constant
action), trajectory-conditioned (unmappable h→action),
trajectory-additive (noise corrupts delta). 800b is the unique working
selector under current constraints. See kills/800b-variants\_step937.md,
Section 4.14.

\subsubsection{7.7 Where the search
points}\label{where-the-search-points}

\textbf{Central question:} Can a substrate accumulate transferable
knowledge without tracking where it has been?

\textbf{Resolved:} (a) D-only substrate navigates at 2x random (800b).
(b) D(s) produces positive R3\_cf (5/7 PASS). (d) Prediction-error
attention achieves R3 encoding (alpha, CONFIRMED). (e) W is sufficient
for encoding but useless for action selection. (h) Epistemic action
selection FAILS (Steps 934-938e). \textbf{(c) Action selection is CLOSED
at 800b under current constraints (Section 4.14).}

\textbf{Open:} (g) Interior point of architecture triangle: encoding
achieved, action selection frozen at 800b. (i) FT09 solvable under
current constraints?

\textbf{(j) Hebbian RNN family (Steps 948-953, DEAD).} First non-916
navigation (seed 8 = 96 L1) but structurally brittle (1/10 seeds, 6
experiments). Requires lucky initialization path. Algorithm invariance
QUALIFIED: non-argmin navigation exists but is unreliable. →
kills/hebbian-rnn\_step953.md

\textbf{(k) Meta-observation: action selection is the bottleneck.}
Post-ban families (916-augmentation ×15, Hebbian RNN ×6,
prediction-selectors ×10, GFS ×3, obs-preprocessing ×2) all fail to
improve action selection beyond 800b. Encoding self-modification works
(alpha, CONFIRMED). But every attempt to learn or modify action
selection breaks navigation. The frozen frame of navigation IS action
selection. This constrains the next family: either accept 800b as fixed
and modify everything else, or find a fundamentally different action
mechanism (ESN, population-based, or architectures that don't yet
exist).

\textbf{(l) Architecture Irrelevance (Proposition 29, Step 955).} ESN
with fixed random \(W_h\) produces numerically identical results to
Hebbian RNN --- same seed 8 = 96 L1, same 9/10 zeros. Architecture
contributes nothing; per-seed variance is entirely determined by the
action RNG's early sequence. →
propositions/29\_architecture\_irrelevance.md

\textbf{(m) The Positive Lock (Proposition 30).} Sigmoid
\(h \in [0,1]^d\) → all dot products positive → first action always
reinforced → winner-take-all after one update. 10\% bootstrap rate from
epsilon-random escape. Lock scales with action space: LS20 (\(n=4\)) ≈
10\%, FT09 (\(n=68\)) ≈ 0\%. Fix: sparse gating (relu threshold) makes
representations state-dependent. → propositions/30\_positive\_lock.md

\textbf{(n) The Temporal Credit Assignment Wall (Proposition 31, Steps
948-990, 43 experiments).} Three interlocking results: (1) 800b's
delta\_per\_action is FROZEN --- 25 modifications all kill LS20. (2)
FT09/VC33 mechanism-limited at any budget (0/10 at 10K, 25K, 50K). (3)
The wall is temporal credit, not the graph ban (Debate v3) ---
parametric state-conditioned models don't solve it. See Section 4.17 for
the formal impossibility result and →
propositions/31\_temporal\_credit\_wall.md.

\textbf{Feasible region = Step 965:} LS20 = 67.0 (chain), FT09 = 0, VC33
= 0, CIFAR = chance. 43 experiments prove this fixed point. Props 29-31
fully characterize the dynamics vertex of the architecture triangle
(Prop 22). The next breakthrough requires a mechanism class outside
prediction-error exploration.

\textbf{(o) Debate v3 Sprint: Simplicity and ARC-AGI-3 (Proposition 32,
Steps 1082-1094, 12 experiments).} Two architecturally distinct substrates
tested on ARC-AGI-3 games with action efficiency scoring (efficiency vs
human baseline, squared). Key findings:

\begin{enumerate}
\item \textbf{Both sides solve 1/3 ARC games at L1 (100\% seeds).} Defense
(reactive switching, zero parameters): ARC=0.30. Prosecution (forward model
action selection in alpha-weighted space): ARC=0.005. Defense 66$\times$
more action-efficient.

\item \textbf{Proposition 32 (Simplicity in R2-constrained learning):}
In the R2-constrained setting (no optimizer), adding learnable parameters
to a working architecture degrades performance. Observed across 2
architecturally distinct families (n=5 experiments: defense v22-v25 killed,
prosecution v21 killed). \emph{Mechanism:} Without SGD, parameter updates
are driven by noisy prediction error via outer product rules. Each additional
parameter accumulates noise at rate $O(\sqrt{T} \cdot \|\mathbf{n}\|)$ vs
signal at rate $O(T \cdot \|\mathbf{s}\|)$, but with $\text{SNR} \ll 1$ in
high-dimensional observation spaces, accumulated noise dominates action
selection before signal reaches actionable levels. Zero-parameter reactive
comparison avoids this entirely. \emph{Prior work:} Simplicity bias in
neural networks is well-studied (Shah et al., NeurIPS 2020; Huh et al.,
2024), but always with gradient-based optimizers that can partially
compensate. No prior work addresses simplicity in the no-optimizer (R2)
setting. \emph{Status:} PROVISIONAL (n=5, post-hoc explanation untested).

\item \textbf{Forward model action selection (prosecution, Steps 1092-1094):}
$W_{\text{fwd}}$ (256$\times$263) predicts next encoding per action in
alpha-weighted space. Action = $\arg\max_a \|W_{\text{fwd}}
[\text{enc}_\alpha, \mathbf{1}_a] - \text{enc}_\alpha\|$. Extended warmup
(100 random steps) required for reliable W\_fwd training (3/10 $\to$ 10/10
seeds). This is R1-compliant ICM (Pathak et al., 2017) without neural
networks or SGD.

\item \textbf{0\% wall persists on $\sim$2/3 games.} Mode 1 (near-inert) =
noise. Mode 2 (responsive, oscillating) = the open problem. Both
architectures fail identically on Mode 2 games.
\end{enumerate}

\subsection{Author Attribution and
Disclosure}\label{author-attribution-and-disclosure}

This research was conducted using a team of LLM agents (Claude Opus and
Claude Sonnet) coordinated by the author. Experiments were designed,
theory was formalized, and the paper was written with LLM assistance.
Experiment scripts were implemented and executed via LLM agents.
Strategic direction, the constitutional framework (R1-R6), approval
gates, and evaluation of findings for self-deception were provided by
the author.

The adversary process (review of each experiment) was conducted by the
LLM agents, not by independent human reviewers. The simulated NeurIPS
review in Section 7 was generated to stress-test the paper's claims. The
constitutional framework (R1-R6) was stress-tested via structured
adversarial debate between two LLM agents (Leo as attacker, Eli as
defender, 12 rounds across two versions). The debate produced: R4's
operational meaning as discriminative capacity (Ashby requisite
variety), R2's role in preventing evaluation hacking (DGM case study),
and Proposition 28 (R3-minimal anti-concentration via Ashby
ultrastability). Full debate log: CONSTITUTIONAL\_DEBATE.md. All
experimental code is available in the repository for independent
verification.

The agents operated on a single machine (Windows 11, RTX 4090) with
experiments run via WSL2. The memory system, mailbox, and coordination
infrastructure are documented in the repository.

\subsection{References}\label{references}

\begin{itemize}
\tightlist
\item
  Abraham, W. C. \& Robins, A. (2005). Memory retention --- the synaptic
  stability versus plasticity dilemma. Trends in Neurosciences, 28(2),
  73-78.
\item
  Bellemare, M. et al.~(2016). Unifying Count-Based Exploration and
  Intrinsic Motivation. NeurIPS.
\item
  Burda, Y. et al.~(2018). Exploration by Random Network Distillation.
  arXiv:1810.12894.
\item
  Burda, Y. et al.~(2019). Large-Scale Study of Curiosity-Driven
  Learning. ICLR.
\item
  Feldman, H. \& Friston, K. J. (2010). Attention, Uncertainty, and
  Free-Energy. Frontiers in Human Neuroscience, 4, 215.
\item
  Da Costa, L. et al.~(2020). Active Inference on Discrete State-Spaces:
  A Synthesis. Journal of Mathematical Psychology, 99, 102447.
\item
  Friston, K. J. (2009). The Free-Energy Principle: A Unified Brain
  Theory? Nature Reviews Neuroscience, 11(2), 127-138.
\item
  Friston, K. J. et al.~(2017). Active Inference: A Process Theory.
  Neural Computation, 29(1), 1-49.
\item
  Fritzke, B. (1995). A Growing Neural Gas Network Learns Topologies.
  NeurIPS.
\item
  Huang, G.-B., Zhu, Q.-Y. \& Siew, C.-K. (2006). Extreme Learning
  Machine: Theory and Applications. Neurocomputing, 70(1-3), 489-501.
\item
  Guo, Z. et al.~(2022). BYOL-Explore: Exploration by Bootstrapped
  Prediction. arXiv:2206.08332.
\item
  Givan, R., Dean, T. \& Greig, M. (2003). Equivalence Notions and Model
  Minimization in Markov Decision Processes. Artificial Intelligence,
  147(1-2), 163-223.
\item
  Graves, A. et al.~(2014). Neural Turing Machines. arXiv:1410.5401.
\item
  Jin, C. et al.~(2020). Reward-Free Exploration for Reinforcement
  Learning. ICML.
\item
  Kirsch, L. \& Schmidhuber, J. (2022). Self-Referential Meta Learning.
  ICML.
\item
  Kakade, S. (2003). On the Sample Complexity of Reinforcement Learning.
  PhD thesis, University College London.
\item
  Kohonen, T. (1988). Self-Organization and Associative Memory.
  Springer.
\item
  Lopes, M., Lang, T. \& Toussaint, M. (2012). Exploration in
  Model-based Reinforcement Learning by Empirically Estimating Learning
  Progress. NeurIPS.
\item
  Oudeyer, P.-Y., Kaplan, F. \& Hafner, V. (2007). Intrinsic Motivation
  Systems for Autonomous Mental Development. IEEE Trans. Evolutionary
  Computation, 11(2).
\item
  Maturana, H. \& Varela, F. (1972). Autopoiesis and Cognition: The
  Realization of the Living.
\item
  McCloskey, M. \& Cohen, N. J. (1989). Catastrophic interference in
  connectionist networks. Psychology of Learning and Motivation, 24,
  109-165.
\item
  Pathak, D. et al.~(2017). Curiosity-driven Exploration by
  Self-Supervised Prediction. ICML.
\item
  Ravindran, B. \& Barto, A. G. (2004). Approximate Homomorphisms: A
  Framework for Non-Exact Minimization in Markov Decision Processes.
  ICML Workshop.
\item
  Rosenstein, M. et al.~(2005). To Transfer or Not To Transfer. NIPS
  Workshop on Inductive Transfer.
\item
  Rudakov, E., Shock, J. \& Cowley, B. U. (2025). Graph-Based
  Exploration for ARC-AGI-3 Interactive Reasoning Tasks.
  arXiv:2512.24156.
\item
  Sajid, N. et al.~(2021). Active Inference: Demystified and Compared.
  Neural Computation, 33(3), 674-712.
\item
  Schmidhuber, J. (1991). Curious Model-Building Control Systems. IEEE
  Int. Joint Conf. on Neural Networks.
\item
  Schmidhuber, J. (2003). Gödel Machines: Self-Referential Universal
  Problem Solvers Making Provably Optimal Self-Improvements.
  arXiv:cs/0309048.
\item
  Schmidhuber, J. (2010). Formal Theory of Creativity, Fun, and
  Intrinsic Motivation. IEEE Trans. Autonomous Mental Development, 2(3).
\item
  Strehl, A. L. \& Littman, M. L. (2008). An Analysis of Model-Based
  Interval Estimation for Markov Decision Processes. JCSS, 74(8),
  1309-1331.
\item
  Tang, H. et al.~(2017). \#Exploration: A Study of Count-Based
  Exploration for Deep RL. NeurIPS.
\item
  van de Ven, G. M. \& Tolias, A. S. (2024). Continual Learning and
  Catastrophic Forgetting. arXiv:2403.05175.
\item
  Wang, Z. et al.~(2019). Characterizing and Avoiding Negative Transfer.
  CVPR.
\item
  Trappe, R. (2026). Phasor Agents: Oscillatory Graphs with Three-Factor
  Plasticity and Sleep-Staged Learning. arXiv:2601.04362.
\item
  Ferret, J. et al.~(2023). A Survey of Temporal Credit Assignment in
  Deep Reinforcement Learning. arXiv:2312.01072.
\item
  Klyubin, A. S., Polani, D. \& Nehaniv, C. L. (2005). Empowerment: A
  Universal Agent-Centric Measure of Control. IEEE CEC.
\item
  Zenil, H. (2026). On the Limits of Self-Improving in Large Language
  Models: The Singularity Is Not Near Without Symbolic Model Synthesis.
  arXiv:2601.05280.
\item
  Kauffman, L. H. (2023). Autopoiesis and Eigenform. Computation,
  11(12), 247.
\item
  Tero, A. et al.~(2010). Rules for Biologically Inspired Adaptive
  Network Design. Science, 327(5964), 439-442.
\item
  Rosen, R. (1991). Life Itself: A Comprehensive Inquiry into the
  Nature, Origin, and Fabrication of Life. Columbia University Press.
\item
  Ashby, W. R. (1956). An Introduction to Cybernetics. Chapman \& Hall.
  {[}Ch. 11: Requisite Variety{]}
\item
  Ashby, W. R. (1960). Design for a Brain. Chapman \& Hall. {[}Ch. 7:
  Ultrastability{]}
\item
  Beer, R. D. (1995). A Dynamical Systems Perspective on
  Agent-Environment Interaction. Artificial Intelligence, 72(1-2),
  173-215.
\item
  Dunning, D. \& Kruger, J. (1999). Unskilled and Unaware of It. Journal
  of Personality and Social Psychology, 77(6), 1121-1134.
\item
  Gould, S. J. \& Eldredge, N. (1972). Punctuated Equilibria. In Schopf,
  T. J. M. (ed.), Models in Paleobiology, 82-115.
\item
  Hernández-Orallo, J. (2017). The Measure of All Minds. Cambridge
  University Press.
\item
  Hubinger, E. et al.~(2019). Risks from Learned Optimization in
  Advanced Machine Learning Systems. arXiv:1906.01820.
\item
  Krueger, D. et al.~(2020). Hidden Incentives and the Design of
  Self-Adaptive Systems. arXiv:2009.09153.
\item
  Lakatos, I. (1978). The Methodology of Scientific Research Programmes.
  Cambridge University Press.
\item
  Li, M. \& Vitányi, P. (2008). An Introduction to Kolmogorov Complexity
  and Its Applications. 3rd ed.~Springer.
\item
  Manheim, D. \& Garrabrant, S. (2018). Categorizing Variants of
  Goodhart's Law. arXiv:1803.04585.
\item
  Popper, K. (1934). The Logic of Scientific Discovery. Routledge.
\item
  Raup, D. M. (1991). Extinction: Bad Genes or Bad Luck? W. W. Norton.
\item
  Rice, H. G. (1953). Classes of Recursively Enumerable Sets and Their
  Decision Problems. Transactions of the AMS, 74(2), 358-366.
\item
  Michaels, J. A. \& Scherberger, H. (2016). hebbRNN: A Reward-Modulated
  Hebbian Learning Rule for Recurrent Neural Networks. JOSS, 1(8), 60.
\item
  Najarro, E. \& Risi, S. (2020). Meta-Learning through Hebbian
  Plasticity in Random Networks. NeurIPS.
\item
  Patel, D. et al.~(2022). Exploration in neo-Hebbian reinforcement
  learning: Computational approaches to the exploration-exploitation
  balance with bio-inspired neural networks. Neural Networks, 151,
  204-218.
\item
  Salge, C., Glackin, C. \& Polani, D. (2014). Empowerment --- An
  Introduction. arXiv:1310.1863.
\item
  Sakana AI (2025). Darwin Gödel Machine: Open-Ended Evolution of
  Self-Improving Agents. arXiv:2505.22954.
\item
  Schlesinger, M. (2014). Learnability Frontier in Learning Agent
  Systems. Cognitive Systems Research, 29-30, 1-18.
\item
  Sutton, R. S. (1992). Adapting Bias by Gradient Descent: An
  Incremental Version of Delta-Bar-Delta. AAAI.
\item
  Wright, S. (1932). The Roles of Mutation, Inbreeding, Crossbreeding
  and Selection in Evolution. Proc. VI Int. Congress of Genetics, 1,
  356-366.
\end{itemize}

\end{document}
